% Options for packages loaded elsewhere
\PassOptionsToPackage{unicode}{hyperref}
\PassOptionsToPackage{hyphens}{url}
\documentclass[
]{article}
\usepackage{xcolor}
\usepackage[margin=1in]{geometry}
\usepackage{amsmath,amssymb}
\setcounter{secnumdepth}{-\maxdimen} % remove section numbering
\usepackage{iftex}
\ifPDFTeX
  \usepackage[T1]{fontenc}
  \usepackage[utf8]{inputenc}
  \usepackage{textcomp} % provide euro and other symbols
\else % if luatex or xetex
  \usepackage{unicode-math} % this also loads fontspec
  \defaultfontfeatures{Scale=MatchLowercase}
  \defaultfontfeatures[\rmfamily]{Ligatures=TeX,Scale=1}
\fi
\usepackage{lmodern}
\ifPDFTeX\else
  % xetex/luatex font selection
\fi
% Use upquote if available, for straight quotes in verbatim environments
\IfFileExists{upquote.sty}{\usepackage{upquote}}{}
\IfFileExists{microtype.sty}{% use microtype if available
  \usepackage[]{microtype}
  \UseMicrotypeSet[protrusion]{basicmath} % disable protrusion for tt fonts
}{}
\makeatletter
\@ifundefined{KOMAClassName}{% if non-KOMA class
  \IfFileExists{parskip.sty}{%
    \usepackage{parskip}
  }{% else
    \setlength{\parindent}{0pt}
    \setlength{\parskip}{6pt plus 2pt minus 1pt}}
}{% if KOMA class
  \KOMAoptions{parskip=half}}
\makeatother
\usepackage{color}
\usepackage{fancyvrb}
\newcommand{\VerbBar}{|}
\newcommand{\VERB}{\Verb[commandchars=\\\{\}]}
\DefineVerbatimEnvironment{Highlighting}{Verbatim}{commandchars=\\\{\}}
% Add ',fontsize=\small' for more characters per line
\usepackage{framed}
\definecolor{shadecolor}{RGB}{248,248,248}
\newenvironment{Shaded}{\begin{snugshade}}{\end{snugshade}}
\newcommand{\AlertTok}[1]{\textcolor[rgb]{0.94,0.16,0.16}{#1}}
\newcommand{\AnnotationTok}[1]{\textcolor[rgb]{0.56,0.35,0.01}{\textbf{\textit{#1}}}}
\newcommand{\AttributeTok}[1]{\textcolor[rgb]{0.13,0.29,0.53}{#1}}
\newcommand{\BaseNTok}[1]{\textcolor[rgb]{0.00,0.00,0.81}{#1}}
\newcommand{\BuiltInTok}[1]{#1}
\newcommand{\CharTok}[1]{\textcolor[rgb]{0.31,0.60,0.02}{#1}}
\newcommand{\CommentTok}[1]{\textcolor[rgb]{0.56,0.35,0.01}{\textit{#1}}}
\newcommand{\CommentVarTok}[1]{\textcolor[rgb]{0.56,0.35,0.01}{\textbf{\textit{#1}}}}
\newcommand{\ConstantTok}[1]{\textcolor[rgb]{0.56,0.35,0.01}{#1}}
\newcommand{\ControlFlowTok}[1]{\textcolor[rgb]{0.13,0.29,0.53}{\textbf{#1}}}
\newcommand{\DataTypeTok}[1]{\textcolor[rgb]{0.13,0.29,0.53}{#1}}
\newcommand{\DecValTok}[1]{\textcolor[rgb]{0.00,0.00,0.81}{#1}}
\newcommand{\DocumentationTok}[1]{\textcolor[rgb]{0.56,0.35,0.01}{\textbf{\textit{#1}}}}
\newcommand{\ErrorTok}[1]{\textcolor[rgb]{0.64,0.00,0.00}{\textbf{#1}}}
\newcommand{\ExtensionTok}[1]{#1}
\newcommand{\FloatTok}[1]{\textcolor[rgb]{0.00,0.00,0.81}{#1}}
\newcommand{\FunctionTok}[1]{\textcolor[rgb]{0.13,0.29,0.53}{\textbf{#1}}}
\newcommand{\ImportTok}[1]{#1}
\newcommand{\InformationTok}[1]{\textcolor[rgb]{0.56,0.35,0.01}{\textbf{\textit{#1}}}}
\newcommand{\KeywordTok}[1]{\textcolor[rgb]{0.13,0.29,0.53}{\textbf{#1}}}
\newcommand{\NormalTok}[1]{#1}
\newcommand{\OperatorTok}[1]{\textcolor[rgb]{0.81,0.36,0.00}{\textbf{#1}}}
\newcommand{\OtherTok}[1]{\textcolor[rgb]{0.56,0.35,0.01}{#1}}
\newcommand{\PreprocessorTok}[1]{\textcolor[rgb]{0.56,0.35,0.01}{\textit{#1}}}
\newcommand{\RegionMarkerTok}[1]{#1}
\newcommand{\SpecialCharTok}[1]{\textcolor[rgb]{0.81,0.36,0.00}{\textbf{#1}}}
\newcommand{\SpecialStringTok}[1]{\textcolor[rgb]{0.31,0.60,0.02}{#1}}
\newcommand{\StringTok}[1]{\textcolor[rgb]{0.31,0.60,0.02}{#1}}
\newcommand{\VariableTok}[1]{\textcolor[rgb]{0.00,0.00,0.00}{#1}}
\newcommand{\VerbatimStringTok}[1]{\textcolor[rgb]{0.31,0.60,0.02}{#1}}
\newcommand{\WarningTok}[1]{\textcolor[rgb]{0.56,0.35,0.01}{\textbf{\textit{#1}}}}
\usepackage{graphicx}
\makeatletter
\newsavebox\pandoc@box
\newcommand*\pandocbounded[1]{% scales image to fit in text height/width
  \sbox\pandoc@box{#1}%
  \Gscale@div\@tempa{\textheight}{\dimexpr\ht\pandoc@box+\dp\pandoc@box\relax}%
  \Gscale@div\@tempb{\linewidth}{\wd\pandoc@box}%
  \ifdim\@tempb\p@<\@tempa\p@\let\@tempa\@tempb\fi% select the smaller of both
  \ifdim\@tempa\p@<\p@\scalebox{\@tempa}{\usebox\pandoc@box}%
  \else\usebox{\pandoc@box}%
  \fi%
}
% Set default figure placement to htbp
\def\fps@figure{htbp}
\makeatother
\setlength{\emergencystretch}{3em} % prevent overfull lines
\providecommand{\tightlist}{%
  \setlength{\itemsep}{0pt}\setlength{\parskip}{0pt}}
\usepackage{bookmark}
\IfFileExists{xurl.sty}{\usepackage{xurl}}{} % add URL line breaks if available
\urlstyle{same}
\hypersetup{
  pdftitle={Fundamentals of ML\_Assignment 1\_Andrew Claxton},
  hidelinks,
  pdfcreator={LaTeX via pandoc}}

\title{Fundamentals of ML\_Assignment 1\_Andrew Claxton}
\author{}
\date{\vspace{-2.5em}2026-05-28}

\begin{document}
\maketitle

Data source:
\url{https://data.wa.gov/demographics/SADE-County-Population-by-Age-and-Sex/8ida-3myk/about_data}

\begin{Shaded}
\begin{Highlighting}[]
\NormalTok{df }\OtherTok{\textless{}{-}} \FunctionTok{read.csv}\NormalTok{(}\StringTok{"C:}\SpecialCharTok{\textbackslash{}\textbackslash{}}\StringTok{Users}\SpecialCharTok{\textbackslash{}\textbackslash{}}\StringTok{User}\SpecialCharTok{\textbackslash{}\textbackslash{}}\StringTok{OneDrive}\SpecialCharTok{\textbackslash{}\textbackslash{}}\StringTok{Desktop}\SpecialCharTok{\textbackslash{}\textbackslash{}}\StringTok{Grad School\_Business Analytics}\SpecialCharTok{\textbackslash{}\textbackslash{}}\StringTok{Fundamentals of ML}\SpecialCharTok{\textbackslash{}\textbackslash{}}\StringTok{Assignment 1}\SpecialCharTok{\textbackslash{}\textbackslash{}}\StringTok{SADE\_County\_Population\_by\_Age\_and\_Sex\_20260528.csv"}\NormalTok{)}

\CommentTok{\#read the CSV}

\FunctionTok{summary}\NormalTok{(df[, }\FunctionTok{c}\NormalTok{(}\StringTok{"County"}\NormalTok{, }\StringTok{"Sex"}\NormalTok{, }\StringTok{"Total.Population.2024"}\NormalTok{, }

               \StringTok{"Total.Population.2023"}\NormalTok{, }\StringTok{"Total.Population.2025"}\NormalTok{)])}
\end{Highlighting}
\end{Shaded}

\begin{verbatim}
##     County              Sex            Total.Population.2024
##  Length:1404        Length:1404        Min.   :    24.59    
##  Class :character   Class :character   1st Qu.:   574.70    
##  Mode  :character   Mode  :character   Median :  1738.60    
##                                        Mean   :  5723.43    
##                                        3rd Qu.:  4217.54    
##                                        Max.   :112750.34    
##  Total.Population.2023 Total.Population.2025
##  Min.   :    23.71     Min.   :    24.06    
##  1st Qu.:   570.71     1st Qu.:   579.99    
##  Median :  1716.61     Median :  1760.77    
##  Mean   :  5663.21     Mean   :  5779.98    
##  3rd Qu.:  4164.82     3rd Qu.:  4293.80    
##  Max.   :110838.03     Max.   :113688.38
\end{verbatim}

\begin{Shaded}
\begin{Highlighting}[]
\CommentTok{\#Print out descriptive statistics for a selection of the data}
\end{Highlighting}
\end{Shaded}

\#This graph shows descriptive statistics for 5 of the 11 variables in
data that captures population figures for counties in Washington State.
The county is a character field containing the names of the counties.
Sex contains either ``Male'' or ``Female''. Then there are three
numerical columns for 2023, 2024, and 2025 population counts.

\begin{Shaded}
\begin{Highlighting}[]
\NormalTok{df}\SpecialCharTok{$}\NormalTok{IsMale }\OtherTok{\textless{}{-}}\NormalTok{ df}\SpecialCharTok{$}\NormalTok{Sex }\SpecialCharTok{==} \StringTok{"Male"} 

\CommentTok{\#Transform Sex column from character to boolean }

\FunctionTok{summary}\NormalTok{(df[, }\FunctionTok{c}\NormalTok{(}\StringTok{"County"}\NormalTok{, }\StringTok{"IsMale"}\NormalTok{, }\StringTok{"Total.Population.2024"}\NormalTok{, }\StringTok{"Total.Population.2023"}\NormalTok{, }\StringTok{"Total.Population.2025"}\NormalTok{)]) }
\end{Highlighting}
\end{Shaded}

\begin{verbatim}
##     County            IsMale        Total.Population.2024 Total.Population.2023
##  Length:1404        Mode :logical   Min.   :    24.59     Min.   :    23.71    
##  Class :character   FALSE:702       1st Qu.:   574.70     1st Qu.:   570.71    
##  Mode  :character   TRUE :702       Median :  1738.60     Median :  1716.61    
##                                     Mean   :  5723.43     Mean   :  5663.21    
##                                     3rd Qu.:  4217.54     3rd Qu.:  4164.82    
##                                     Max.   :112750.34     Max.   :110838.03    
##  Total.Population.2025
##  Min.   :    24.06    
##  1st Qu.:   579.99    
##  Median :  1760.77    
##  Mean   :  5779.98    
##  3rd Qu.:  4293.80    
##  Max.   :113688.38
\end{verbatim}

\begin{Shaded}
\begin{Highlighting}[]
\CommentTok{\#Reprint descriptive statistics to check corrected column data types}
\end{Highlighting}
\end{Shaded}

\#This is the same descriptive statistics, but now we have transformed
the Sex column into a Boolean called ``IsMale'' containing an array of
values that are equal to either TRUE or FALSE

\begin{Shaded}
\begin{Highlighting}[]
\FunctionTok{plot}\NormalTok{(}\FunctionTok{density}\NormalTok{(df}\SpecialCharTok{$}\NormalTok{Total.Population}\FloatTok{.2025}\NormalTok{), }\AttributeTok{main =} \StringTok{"Density of Total Population 2025"}\NormalTok{, }\AttributeTok{xlab =} \StringTok{"Population"}\NormalTok{) }
\end{Highlighting}
\end{Shaded}

\pandocbounded{\includegraphics[keepaspectratio]{Assignment-1_Claxton_files/Assignment-1_Claxton_files/figure-latex/unnamed-chunk-3-1.pdf}}

\begin{Shaded}
\begin{Highlighting}[]
\CommentTok{\#Plot at least one quantitative variable. In this case, a density plot of the Total.Population.2025 column}
\end{Highlighting}
\end{Shaded}

\#This shows the density plot of values from the column containing total
population from 2025. As you can easily see, there is a large amount of
lower values and a long tail to the right. This is normal for population
figures, where most counties will have a smaller number of occupants but
a few will have much larger numbers of occupants.

\begin{Shaded}
\begin{Highlighting}[]
\NormalTok{adams\_df }\OtherTok{\textless{}{-}} \FunctionTok{subset}\NormalTok{(df, County }\SpecialCharTok{==} \StringTok{"Adams"}\NormalTok{) }

\CommentTok{\#create subset of rows with only Adams County }

\NormalTok{adams\_df}\SpecialCharTok{$}\NormalTok{AgeRanges }\OtherTok{\textless{}{-}} \DecValTok{5} \SpecialCharTok{*}\NormalTok{ adams\_df}\SpecialCharTok{$}\NormalTok{Age.Sorting.Variable}\SpecialCharTok{{-}} \DecValTok{3} 

\CommentTok{\#create new numeric column to turn the age sorting variable into a regular age (slightly reduces granularity of data, but useful for plotting) }

\FunctionTok{plot}\NormalTok{(adams\_df}\SpecialCharTok{$}\NormalTok{AgeRanges, adams\_df}\SpecialCharTok{$}\NormalTok{Total.Population}\FloatTok{.2025}\NormalTok{, }\AttributeTok{xlab =} \StringTok{"Age (0{-}85+ years)"}\NormalTok{, }\AttributeTok{ylab =} \StringTok{"2025 Population"}\NormalTok{, }\AttributeTok{main =} \StringTok{"Adams County 2025 Population by Age"}\NormalTok{) }
\end{Highlighting}
\end{Shaded}

\pandocbounded{\includegraphics[keepaspectratio]{Assignment-1_Claxton_files/Assignment-1_Claxton_files/figure-latex/unnamed-chunk-4-1.pdf}}

\begin{Shaded}
\begin{Highlighting}[]
\CommentTok{\#plot scatter chart of Adams County 2025 population by Age }
\end{Highlighting}
\end{Shaded}

\#This is a basic scatter plot showing the 2025 population figures by
age group for only Adams county. The original data had age brackets with
an age sorting variable to track ages groups (0-4, 5-9, 9-13, etc.). We
converted this to a numerical value by using the middle value of those
brackets, which loses some of the data's granularity but is useful for
visualization

\begin{Shaded}
\begin{Highlighting}[]
\NormalTok{colorForGraph }\OtherTok{=} \FunctionTok{ifelse}\NormalTok{(df}\SpecialCharTok{$}\NormalTok{IsMale}\SpecialCharTok{==}\ConstantTok{TRUE}\NormalTok{, }\StringTok{"blue"}\NormalTok{, }\StringTok{"red"}\NormalTok{)}
\CommentTok{\#establish variable for color coding male and female data points}

\FunctionTok{plot}\NormalTok{(adams\_df}\SpecialCharTok{$}\NormalTok{AgeRanges, adams\_df}\SpecialCharTok{$}\NormalTok{Total.Population}\FloatTok{.2025}\NormalTok{, }\AttributeTok{xlab =} \StringTok{"Age (0{-}85+ years)"}\NormalTok{, }\AttributeTok{ylab =} \StringTok{"2025 Population"}\NormalTok{, }\AttributeTok{main =} \StringTok{"Adams County 2025 Population by Age"}\NormalTok{, }\AttributeTok{col=}\NormalTok{colorForGraph) }

\FunctionTok{legend}\NormalTok{(}\StringTok{"topright"}\NormalTok{,}
       \AttributeTok{legend =} \FunctionTok{c}\NormalTok{(}\StringTok{"Male"}\NormalTok{, }\StringTok{"Female"}\NormalTok{),}
       \AttributeTok{col =} \FunctionTok{c}\NormalTok{(}\StringTok{"blue"}\NormalTok{, }\StringTok{"red"}\NormalTok{),}
       \AttributeTok{pch =} \DecValTok{16}\NormalTok{)}
\end{Highlighting}
\end{Shaded}

\pandocbounded{\includegraphics[keepaspectratio]{Assignment-1_Claxton_files/Assignment-1_Claxton_files/figure-latex/unnamed-chunk-5-1.pdf}}

\begin{Shaded}
\begin{Highlighting}[]
\CommentTok{\#reprint scatter plot with colors and a legend.}
\end{Highlighting}
\end{Shaded}

\#This is the same scatter plot with males and females color coded to
give better clarity to the population totals. Thank You!

\end{document}
